\begin{center}
    {\large \textbf{Giorno 6} \textbar\ La falsa Lombok \textbar\ 14 Agosto 2024 \textbar\ \textbf{Kuta}, Lombok}
\end{center}

\noindent\rule{\textwidth}{0.4pt}
\begin{figure}[h!] % Portrait images below
    \centering
    % First portrait image (on the left)
    \begin{minipage}[h]{0.48\textwidth} % Set width equal to half the page
        \centering
        \includegraphics[width=\textwidth]{images/flags/Screenshot Kuta Lombok.png}
    \end{minipage}
    \hfill
    % Text
    \begin{minipage}[h]{0.48\textwidth}
        \raggedright
        \textbf{Superficie}: 4.567 km² \\
        \vspace{0.3cm}
        \textbf{Popolazione}: 3,3 milioni \\
        \vspace{0.3cm}
        \textit{L'isola}
    \end{minipage}
\end{figure}

\landscapeLargeMargin{images/day6/PXL_20240814_000411457.jpg}

\landscapeLargeMargin{images/day6/PXL_20240814_005356814.jpg}

\landscapeLargeMargin{images/day6/PXL_20240814_041156709.jpg}

\landscapeLargeMargin{images/day6/PXL_20240814_041839209.jpg}

\landscapeLargeMargin{images/day6/PXL_20240814_051839737.jpg}

\landscapeLargeMargin{images/day6/PXL_20240814_052120074.jpg}

\landscapeLargeMargin{images/day6/PXL_20240814_052123368.jpg}

\landscapeLargeMargin{images/day6/PXL_20240814_071739934.jpg}

\noindent 14 Agosto 2024 \newline

Anche questa notte sembra procedere tranquilla, fino alle quattro del mattino che ormai sono diventate un appuntamento fisso. Sono svegliato dal canto del gallo, che immagino a sua volta sia stato svegliato dal Muezzin. Abbiamo scoperto infatti che a Lombok si trova una moschea sulla costa e quando è ora di preghiera la sentiamo fino sull'isola, per fortuna abbastanza attutita dal rumore del mare. I galli però non la pensano come noi e sono belli incazzati, e quello che mi tiene sveglio in particolare è vicinissimo alla nostra casetta. Per un attimo mi prendo male e temo che il chicchirichì provenga da Elisa, che dopo aver cenato con una zuppa di mais si è trasformata in gallina. Per fortuna controllo ed è sempre lei, che prova a dormire con un orecchio sul cuscino e l'altro coperto da una mano per isolare il suono esterno. Ai galli si unisce anche un animale che sembra un picchio, che si trova sopra la casa e fa un tac tac tac atroce sulla struttura del tetto. Insomma è il solito circo, e visto che stavolta il gallo sembra vicino vorrei uscire in stile Pozzetto in Ragazzo di campagna, quando lancia uno stivale al povero animale dopo che questo l'ha svegliato al mattino.

Quando mi risveglio alle sette ragiono sulle possibili scuse per evitare il corso mattutino di yoga senza sentirmi in colpa. Viene abbastanza facile perché in effetti non si trattava di yoga bensì di meditazione, e mi immagino già Eli che mi guarda storto mentre l'istruttrice fa "OMMMMMMMM". Convinco facilmente anche lei a desistere.

Ci alziamo mezz'ora dopo per andare a fare colazione e passiamo di fronte alla lezione di meditazione, dove sono effettivamente tutti seduti con le gambe incrociate. La sera precedente mi era capitato di sentire la conversazione di una coppia di italiani nella casetta limitrofa e ho scoperto che l'istruttrice è una ragazza occidentale che dopo aver girato il mondo facendo yoga si è innamorata di questo posto e ha deciso di stabilirsi qui collaborando con il resort. C'è in effetti un foltissimo programma di lezioni private in bella mostra su un tableau vicino alla reception.

Questa mattina per colazione Elisa riprova la English breakfast mentre per me la Continental: uovo sodo, pane tostato con burro d'arachidi e una brioche mignon, il tutto accompagnato da frutta fresca spruzzata di lime e frullato di anguria e papaya. Divoriamo tutto come al solito e mentre ci fanno il conto finale andiamo a prendere i bagagli e a liberare la stanza.

Dopo aver saldato mi chiedono di compilare un questionario sul soggiorno e non posso che dare voti altissimi. Mentre li ringrazio mi avvisano che se volessi tornare o consigliarli ad amici potrebbero farmi un prezzo di favore, ovviamente evitando Booking.com.

La barca che ci deve riportare a Lombok dovrebbe arrivare a momenti e ne approfittiamo per fare due passi e riempirci gli occhi di memorie. Non facciamo in tempo ad arrivare ai lettini che vediamo il personale del resort caricare i nostri bagagli in barca. Corriamo, li ringraziamo nuovamente e, togliendo le scarpe, saliamo sul catamarano. Mentre facciamo foto e video della partenza loro sono ancora sulla spiaggia a salutarci agitando la mano.

Il viaggio è breve come quello dell'andata e nel frattempo il nostro autista mi avvisa di essere già arrivato al punto di ritrovo. Si chiama Dian, ci siamo scritti per l'orario e la logistica dell'appuntamento ed è molto gentile e disponibile. Ci accompagna velocemente al Novotel e lungo il tragitto rallenta nei pressi di dove avevo visto i macachi all'andata, per permettere a Elisa di scattare fotografie.

A un certo punto passiamo davanti a una pista da corsa e scopriamo che qui si svolge il Gran Premio d'Indonesia della MotoGP. Anche Dian, come ogni indonesiano che si rispetti, è un patito di Valentino Rossi.

Quando ci lascia davanti alla reception del Novotel rimaniamo per un attimo diffidenti e delusi. Siamo stati due giorni nella giungla a vivere coi nativi e subito dopo due giorni su una spiaggia tropicale dove le persone sorridevano dalla mattina alla sera. Qui ci ritroviamo catapultati in una location turistica con comfort e intrattenimento occidentale, dove le persone non sono gentili ma gli è imposto di essere servizievoli coi turisti.

Anche alla reception la situazione è diversa: paga subito, carta di credito, documenti, verifiche, scartoffie. Ci consegnano le chiavi e ci accompagnano nella camera: grande e spaziosa, ma pur sempre una camera d'albergo standard, con niente di caratteristico se non il porta-fazzoletti in madreperla. La finestra si affaccia sul circuito di MotoGP e possiamo assistere ai giri di allenamento degli amatori — cosa che per qualcuno sarebbe un plus, ma a me mancano già gli animali e le onde del mare.

Andiamo subito a esplorare il resto dell'hotel, costruito ispirandosi alle case dei nativi antiche ma in maniera del tutto snaturata, e ci dirigiamo verso la spiaggia dove nessuno fa il bagno in mare perché sono tutti in piscina. Abbiamo visto abbastanza.

Andiamo alla reception a chiedere quanto disti Kuta a piedi e dopo aver ragionato sulle possibili alternative decidiamo di farci accompagnare da un autista dell'hotel, che per pochi euro ci sarebbe venuto a prendere anche al ritorno all'ora che desideravamo.

Anche Kuta all'inizio ci lascia molto perplessi: bancarelle ai lati della strada principale con t-shirt e gadget occidentali, surfisti coi capelli lunghi che entrano nei supermercati scalzi e senza maglietta, turisti da tutte le parti e nativi che si dedicano solo alla vendita di tour e souvenir cercando di attirarti a gran voce nei loro stand.

Entriamo in un supermercato che fa anche gastronomia e prendiamo due prodotti freschi: un medaglione di verdure e patate per Eli e un tortino di pollo per me, oltre a due gelati — uno all'avocado e uno all'Oreo. Il medaglione era buono; riguardo al tortino di pollo, non starò tranquillo finché non arriverò sano e salvo a Torino.

Ci sediamo di fronte al supermercato e cerchiamo di capire il da farsi, perché al momento siamo molto scoraggiati. Trovo di nuovo un'analogia col Messico, quando ci siamo ritrovati a Cancún senza auto, disperati per lo scam dei tassisti.

I tour del Novotel non sono neanche da prendere in considerazione. Quando li avevo letti a Torino mi erano sembrati interessanti e ragionevoli, ma dopo qualche giorno qui la prospettiva è cambiata: il tour delle risaie non avrebbe potuto battere quello del Tetebatu fatto con Oshi, il tour di snorkeling non avrebbe potuto superare la spiaggia di coralli di Gili Asahan, di fare trekking non abbiamo particolarmente voglia e i due tour rimanenti sono presso spiagge affollate.

Decidiamo allora di girare per Kuta a contrattare un tour autentico nell'entroterra, a costo di farcelo disegnare su misura. Scopriamo che Kuta ha due vie principali: quella turistica su cui ci siamo avventurati finora e una strada che la taglia perpendicolarmente, frequentata principalmente dalla popolazione locale. È qui che ci dirigiamo ed è qui che iniziamo a cercare un tour.

Al primo posto a cui chiediamo ci dicono di tornare dopo perché il proprietario sta pranzando. Mi piace già come si ragiona qui: “chill, relax first”. Al secondo posto interagiamo con un ragazzo simpatico, ma all'apparenza uno che vuole fregare i turisti. Lo teniamo venti minuti occupato a spiegarci le varie alternative per una giornata intera. Stiamo quasi per concordare una giornata in barca a pescare e mangiare pesce, quando decidiamo di tornare al piano originale: un tour guidato per i villaggi locali.

La contrattazione sembra quasi una negoziazione di ostaggi. L'accordo finale è per un autista privato che ci porti a vedere tre villaggi e infine la spiaggia di Tanjung Aan fino al tramonto. Non è facile chiarirsi e ci sono sempre zone d'ombra nell'itinerario — il "fino al tramonto" che cozza con "il tour termina alle quattro" — ma continuando a insistere arriviamo a una quadra. Come precauzione, durante tutta la conversazione lo registro col cellulare. Forse un po' paranoico, ma ci serve anche a tenere traccia dell'itinerario concordato. La cifra finale è assolutamente sensata e non paragonabile a quella dei carissimi tour organizzati dal Novotel.

Contenti di aver sistemato i piani per l'indomani, dobbiamo risolvere la giornata odierna e decidiamo di addentrarci ancora di più nell'interno di Kuta, per la stradina dei locals, in modo da vedere come vivono e riuscire a scattare foto autentiche. Eli ha la Reflex al collo e ovviamente mi torna in mente il Messico, quando a San Cristóbal de las Casas sembravo il bodyguard assoldato da Elisa per proteggerla. Allora anche l'abbigliamento era adeguato, mentre questa volta ho una canotta bianca e blu e un costume da bagno azzurro che arriva a metà coscia. L'altra volta potevo sembrare una guardia del corpo, qui sembro un surfista albino che non riesce a prendersi un'abbronzatura uniforme. Un "coglionazzo" insomma.

Per fortuna tutti ci sorridono e ci salutano e non c'è bisogno di stare sulla difensiva. Entriamo nei negozi dei locali e proviamo qualche loro prodotto: in particolare le chips di banana, che abbiamo imparato ad amare in Messico, e le chips di patata viola dolce, che ci fanno letteralmente impazzire. Proviamo anche le sigarette locali, curiosi dopo i discorsi sul tabacco diverso da quello della Philip Morris, ma l'esperimento non va a buon fine. Queste contengono 2,2 mg di nicotina per sigaretta: per intenderci, le Marlboro rosse — tra quelle con la percentuale più alta qui in Italia — ne hanno 0,8 mg. Facciamo un tiro per provarle e ci viene un mal di testa immediato, che per fortuna passa in fretta. Compriamo comunque del tabacco sfuso da regalare.

Mentre giriamo per i negozietti vediamo finalmente da vicino i gruppi scolastici di ragazzi che marciano e facciamo foto e video a profusione. Continuando nella passeggiata arriviamo fino al termine della città e finalmente riesco a convincere Elisa a tornare indietro.

Qui veniamo quasi inseguiti da tre bambine a cui Elisa ha scattato fotografie: dapprima ci salutano, poi iniziano a chiederci "Money! Money!". Io sono per andarmene senza voltarmi, ma Elisa-maestra è meno cuore di pietra di me e compra alle bambine tre gelati — più un quarto per me, perché se lo mangiano loro allora lo voglio anch'io. Una bomba di cioccolato e praline.

Torniamo al punto di partenza dove il nostro autista è già ad aspettarci dopo che lo abbiamo avvisato su WhatsApp, e ritorniamo in hotel. Proviamo a goderci l'ultima ora di sole in spiaggia, sui lettini del Novotel. Sono comodissimi e l'insenatura è davvero bella, ma non vale il giro per le strade sporche e afose di Kuta. Il sole tramonta in fretta e con il buio scende anche la temperatura, quindi andiamo in camera a rinfrescarci e a prepararci per la cena.

Decidiamo di provare il ristorante del Novotel, dal momento che anche il menù consultato online ci ispirava. Quando arriviamo ci comunicano che quella sera c'è il menù a buffet. Strike 1. Va bene, accettiamo, fiduciosi che sarà migliore di quello dell'hotel a Jakarta. La qualità è effettivamente migliore, ma la scelta lascia abbastanza a desiderare. Strike 2. Quando siamo seduti a tavola scopriamo che, insistendo un po', vari ospiti sono riusciti a ordinare alla carta — possibilità che a noi non avevano esposto. Strike 3. Battitore eliminato.

Ci facciamo coraggio e cerchiamo di combinare gli ingredienti giusti. Il tofu fritto con foglia di spinacio è ottimo, il pesce in pastella non è male anche se potevano evitare di immergerlo nella maionese, c'è una zuppetta tipica molto buona e mi faccio preparare al momento un involtino vietnamita che mi lascia soddisfatto. Dal buffet provo anche ad assaggiare la carbonara. "Stop! Stop! He's already dead!" Ovviamente sa di penne con panna e prosciutto, ma almeno la cottura si salva.

Anche il tentativo di provare una birra nuova sembra fallire, perché ordino erroneamente la loro birra al limone, ma alla fine somiglia a una panaché e il sapore non mi dispiace.

Rimaniamo per un po' sul terrazzino del ristorante vista spiaggia, principalmente a ricordare i noodles di Oshi, e quando inizia a fare un po' più fresco torniamo in camera.

Per far capire quanto non sono abituato a questo genere di location: abbiamo finito l'acqua e scendo alla reception a chiedere dove trovo un distributore senza dover andare fino al ristorante. Mi dicono che me la possono mandare in camera e quando chiedo un paio di bottiglie me ne vogliono mandare quattro. Rispondo un po' seccato "Two is fine" e me ne vado. Scoprirò il giorno dopo che l'acqua è gratuita e che si poteva ordinare tranquillamente chiamando il servizio in camera dalla stanza.

\pagestyle{empty} % disable numbers

%single
\twoImagesLargeMargins{images/day6/IMG_4029.JPG}
                      {images/day6/IMG_4032.JPG}
%

%coppia
\landscapeLargeMargin{images/day6/PXL_20240814_042226531.MP.jpg}
\landscapeLargeMarginCentered{images/day6/PXL_20240814_075142484.jpg}
%

%coppia
\landscapeThinMarginCentered{images/day6/IMG_4163.JPG}
\twoImagesLargeMargins{images/day6/IMG_4186.JPG}
                     {images/day6/IMG_4108.JPG}
%

%coppia
\twoImagesLargeMargins{images/day6/IMG_4239.JPG}
                      {images/day6/IMG_4248.JPG}
\landscapeThinMarginCentered{images/day6/IMG_4219.JPG}
%

%coppia
\landscapeThinMargin{images/day6/IMG_4275.JPG}
\twoImagesLargeMargins{images/day6/IMG_4329.JPG}
                      {images/day6/IMG_4339.JPG}
%

%coppia
\twoImagesThinMargins{images/day6/IMG_4356.JPG}
                     {images/day6/IMG_4367.JPG}
\twoImagesThinMargins{images/day6/IMG_4401.JPG}
                     {images/day6/IMG_4399.JPG}
%

%coppia
\landscapeThinMarginCentered{images/day6/IMG_4410.JPG}
\landscapeThinMarginCentered{images/day6/IMG_4423.JPG}
%

%coppia
\landscapeThinMarginCentered{images/day6/IMG_4380.JPG}
\landscapeThinMarginCentered{images/day6/IMG_4456.JPG}
%

%coppia
\splitImageLeft{images/day6/PXL_20240814_094417639.jpg}
\splitImageRight{images/day6/PXL_20240814_094417639.jpg}
%

\landscapeLargeMarginCentered{images/day6/PXL_20240814_113359359.jpg}

\pagestyle{fancy} % enable numbers